\begin{description}
\item[1行目] ファイルからデータを読み込みます。このデータファイル\verb|baseballDecade.csv|には野手，投手の10年分のデータが入っています。
\item[3-15行目] 最初のブロックとして，データの加工・変改を行い，それを\verb|dat|オブジェクトに格納します。
\begin{description}
	\item[4行目] フィルターをかけて，投手だけのデータに絞り込みます。
	\item[5行目] フィルターをかけて，Swallowsの選手に限ります
	\item[6行目] 選手名でグルーピングします。
	\item[7行目] \verb|nest|関数でグループ化変数に沿ってデータを畳み込みます。引数をとくに指定しなければ，変数はすべて\verb|data|という変数名に含まれます。この段階でデータフレームは出力\ref{RS:out25_01}のようになっています。\verb|nest|関数のイメージを図\ref{fig:25_03}に示しました。
	\item[8行目] 畳み込まれたデータセットを使って，そのデータの行数(何年分のデータがあるか\footnote{ここで\verb|data|とあるのはまとめられたミニ・データフレーム全体を引数にしています。それに対して\verb|NROW|という関数をあてがうことで行数を数えています。\verb|NROW|関数の引数は\verb|.|ですが，これは\verb|map|関数の第一引数全体を再度参照するときの略記号です。\verb|purrr::map|というのは\verb|purrr|パッケージの\verb|map|関数で，同じ関数操作を繰り返し行うときに使います。})，勝利数情報の欠損値の有無を作る変数\footnote{\verb|purrr::map|の使い方は先ほどと同じで，\verb|.$Win|は\verb|data$Win|と同じ意味です。\verb|anyNA|はNAがあるかどうかを\verb|TRUE/FALSE|で返す関数です。}を作っています。
	\item[9行目] フィルターをかけて，データの行数\verb|n|が7より大きいもの，つまり8年以上登板している一流の選手だけに絞ったのです。
	\item[10行目] フィルターをかけて，欠損値がないデータだけに絞っています。
	\item[11行目] 畳み込みを解除します。
	\item[12行目] 年度，選手名，年俸，勝利数だけ変数をセレクトして取り出します。
\end{description}
